% results_and_conclusion.tex
% Sections 4–6 of the SGFood-233 research paper
% All accuracy values verified against notebook outputs and JSON result files.

% -----------------------------------------------------------------
% SECTION 4 — RESULTS
% -----------------------------------------------------------------
\section{Results}\label{sec:results}

\subsection{Baseline Architecture Comparison}\label{sec:baseline}

Table~\ref{tab:baseline} summarises the performance of the three architectures
trained under identical conditions on the original (uncleaned) 75/25
train/validation split (lr$\,{=}\,$1e-4, weight decay$\,{=}\,$1e-4,
batch size$\,{=}\,$64, cosine-annealing schedule, 10 epochs).

\begin{table}[ht]
\centering
\caption{Baseline architecture comparison on the original 75/25 split
         (identical hyperparameters).}
\label{tab:baseline}
\begin{tabular}{lccc}
\toprule
\textbf{Model} & \textbf{Val Accuracy (\%)} & \textbf{Macro F1} & \textbf{Weighted F1} \\
\midrule
ResNet-50   & 75.09 & 0.736 & --- \\
ViT-Base    & 80.16 & 0.790 & --- \\
ConvNeXt-Base & 83.45 & 0.825 & --- \\
\bottomrule
\end{tabular}
\end{table}

ConvNeXt-Base outperforms the CNN baseline (ResNet-50) by
$+8.36$ percentage points (pp) and the Vision Transformer (ViT-Base) by
$+3.29$~pp.  All three models were pre-trained on ImageNet-1K and
fine-tuned with the same augmentation pipeline (random crop to
$224{\times}224$, horizontal flip, ImageNet normalisation).  The
consistent ranking under identical optimisation conditions establishes
ConvNeXt-Base as the strongest single architecture on SGFood-233 under
default hyperparameters.

% -----------------------------------------------------------------
\subsection{Error Analysis}\label{sec:error-analysis}

\subsubsection{Per-Class Performance}\label{sec:per-class}

Despite strong aggregate metrics, per-class performance varies
substantially.  Table~\ref{tab:f1-gap} lists the classes exhibiting the
largest F1 gap between the best and worst model, revealing where
architectural choice matters most.

\begin{table}[ht]
\centering
\caption{Classes with the largest inter-model F1 gap (top~10).}
\label{tab:f1-gap}
\begin{tabular}{lcccc}
\toprule
\textbf{Class} & \textbf{CNN F1} & \textbf{ViT F1} & \textbf{ConvNeXt F1} & \textbf{Gap} \\
\midrule
Don, chicken teriyaki & 0.452 & ---   & 0.662 & 0.210 \\
tortilla -- plain     & 0.372 & ---   & 0.578 & 0.206 \\
Barley                & 0.304 & ---   & 0.504 & 0.200 \\
satay bee hoon        & 0.521 & ---   & 0.715 & 0.194 \\
wholegrain muffin     & 0.671 & ---   & 0.860 & 0.189 \\
roti john             & 0.676 & ---   & 0.860 & 0.184 \\
Mee siam, fried       & 0.273 & ---   & 0.452 & 0.179 \\
Macaroni              & 0.629 & ---   & 0.802 & 0.173 \\
Peanut Pancake        & 0.737 & ---   & 0.909 & 0.172 \\
ngoh hiang            & 0.582 & ---   & 0.754 & 0.172 \\
\bottomrule
\end{tabular}
\end{table}

The classes with the widest gaps are predominantly noodle and
grain-based dishes that require both textural and compositional
understanding.  ConvNeXt closes these gaps consistently, suggesting
its hybrid receptive field is especially beneficial for fine-grained
food categories.

\subsubsection{Confusion Pairs}\label{sec:confusion-pairs}

Table~\ref{tab:confusion} lists the most frequently confused class
pairs for each architecture, drawn from the validation-set confusion
matrices.

\begin{table}[ht]
\centering
\caption{Top-5 confused class pairs per model (validation set,
         original 75/25 split).}
\label{tab:confusion}
\small
\begin{tabular}{llc}
\toprule
\textbf{Model} & \textbf{True $\to$ Predicted} & \textbf{Count} \\
\midrule
\multirow{5}{*}{CNN}
  & bee hoon $\to$ Bee hoon, goreng         & 50 \\
  & char siew pau $\to$ steamed buns        & 48 \\
  & Sweets $\to$ Chocolate                  & 46 \\
  & chicken soup $\to$ vegetable soup       & 46 \\
  & har cheong gai $\to$ fried chicken      & 46 \\
\midrule
\multirow{5}{*}{ViT}
  & Lontong w/ sayur lodeh $\to$ lontong    & 44 \\
  & Sweets $\to$ Chocolate                  & 41 \\
  & char siew pau $\to$ steamed buns        & 41 \\
  & steamed buns $\to$ char siew pau        & 37 \\
  & Seafood Noodles Soup $\to$ sliced fish soup & 35 \\
\midrule
\multirow{5}{*}{ConvNeXt}
  & char siew pau $\to$ steamed buns        & 44 \\
  & Sweets $\to$ Chocolate                  & 43 \\
  & steamed buns $\to$ char siew pau        & 41 \\
  & Lontong w/ sayur lodeh $\to$ lontong    & 38 \\
  & chicken soup $\to$ vegetable soup       & 37 \\
\bottomrule
\end{tabular}
\end{table}

Each architecture exhibits a distinct error signature.
\textbf{CNN} uniquely confuses texture-similar noodle dishes
(\emph{bee hoon}~$\leftrightarrow$~\emph{Bee hoon, goreng}), where the
primary visual difference is sauce colour rather than shape.
\textbf{ViT} confuses structurally similar dishes
(\emph{Lontong with sayur lodeh}~$\leftrightarrow$~\emph{lontong}),
where the distinction lies in accompanying condiments rather than the
main item.
\textbf{ConvNeXt} shares the universal \emph{char siew pau}
$\leftrightarrow$ \emph{steamed buns} confusion---an inherently
ambiguous pair---but produces fewer model-specific failure modes,
consistent with its balanced feature extraction.

% -----------------------------------------------------------------
\subsection{Model Interpretability via Grad-CAM}\label{sec:gradcam}

Gradient-weighted Class Activation Maps (Grad-CAM)~\cite{selvaraju2017grad}
were generated for correctly classified, misclassified, and model-disagreement
samples (Figure~\ref{fig:gradcam}).

\begin{figure}[ht]
\centering
% \includegraphics[width=\textwidth]{figures/gradcam_comparison.pdf}
\textbf{[Placeholder: Grad-CAM comparison grid --- CNN / ViT / ConvNeXt
across correctly classified, CNN-only errors, ViT-only errors, and
all-model errors.]}
\caption{Grad-CAM attention maps for the three architectures on
         representative SGFood-233 samples.}
\label{fig:gradcam}
\end{figure}

Three distinct attention strategies emerge:
\begin{itemize}
  \item \textbf{CNN (ResNet-50):} Activation concentrates on local
    texture patches---colour gradients, surface patterns---explaining
    its strength on visually distinct items (e.g.\ fruits) and weakness
    on texture-similar noodle variants.
  \item \textbf{ViT-Base:} Attention distributes across the global dish
    structure and key components (e.g.\ bowl shape, garnishes),
    capturing compositional relationships but occasionally conflating
    structurally similar presentations.
  \item \textbf{ConvNeXt-Base:} Activation balances local detail and
    global context, attending simultaneously to texture cues and overall
    dish layout.  This complementary feature capture aligns with its
    superior accuracy and fewer systematic confusion patterns.
\end{itemize}

% -----------------------------------------------------------------
\subsection{Computational Efficiency}\label{sec:efficiency}

Table~\ref{tab:efficiency} benchmarks inference-time efficiency on a
single NVIDIA GPU (batch size 1, mixed precision).

\begin{table}[ht]
\centering
\caption{Computational efficiency comparison.}
\label{tab:efficiency}
\begin{tabular}{lcccc}
\toprule
\textbf{Model} & \textbf{Params (M)} & \textbf{FLOPs (G)} &
  \textbf{Throughput (img/s)} & \textbf{Peak Mem (MB)} \\
\midrule
ResNet-50     & 24.0  & 4.1  & 211.2 & 2\,315 \\
ViT-Base      & 86.0  & 16.9 & 142.8 & 2\,218 \\
ConvNeXt-Base & 87.8  & 15.4 & 117.4 & 2\,536 \\
\bottomrule
\end{tabular}
\end{table}

ResNet-50 is 1.8$\times$ faster than ConvNeXt-Base while requiring
3.7$\times$ fewer parameters and 3.8$\times$ fewer FLOPs.  However,
this speed advantage comes at the cost of $-8.36$~pp in accuracy.
ConvNeXt-Base, despite its larger footprint, remains practical for
server-side deployment at 117~images/s.  Interestingly, ViT-Base
consumes the most FLOPs (16.9\,G) yet achieves neither the highest
accuracy nor the highest throughput, placing it off the Pareto frontier
in the accuracy--efficiency trade-off (Figure~\ref{fig:pareto}).

\begin{figure}[ht]
\centering
% \includegraphics[width=0.7\textwidth]{figures/pareto_efficiency.pdf}
\textbf{[Placeholder: Pareto plot --- throughput (img/s) vs.\ validation
accuracy for the three architectures.]}
\caption{Accuracy vs.\ throughput Pareto frontier.}
\label{fig:pareto}
\end{figure}

% -----------------------------------------------------------------
\subsection{Embedding Space Analysis}\label{sec:embeddings}

To evaluate learned representation quality beyond classification
accuracy, we extracted penultimate-layer embeddings for 5\,000
validation samples per model and computed silhouette scores in both
the original high-dimensional space and after t-SNE reduction to 2D
(Table~\ref{tab:silhouette}).

\begin{table}[ht]
\centering
\caption{Silhouette scores measuring class separability in the learned
         embedding space.}
\label{tab:silhouette}
\begin{tabular}{lccc}
\toprule
\textbf{Space} & \textbf{CNN} & \textbf{ViT} & \textbf{ConvNeXt} \\
\midrule
High-dimensional & 0.134 & 0.129 & 0.200 \\
t-SNE (2D)       & 0.310 & 0.426 & 0.502 \\
\bottomrule
\end{tabular}
\end{table}

ConvNeXt-Base produces the most separable embeddings in both spaces,
with a high-dimensional silhouette score 49\% higher than the CNN and
55\% higher than the ViT.  The t-SNE projections
(Figure~\ref{fig:tsne}) visually confirm this: ConvNeXt clusters are
tighter and better separated, while CNN clusters appear more dispersed
with greater inter-class overlap.

\begin{figure}[ht]
\centering
% \includegraphics[width=\textwidth]{figures/tsne_umap_comparison.pdf}
\textbf{[Placeholder: Side-by-side t-SNE and UMAP projections for
CNN, ViT, and ConvNeXt.]}
\caption{t-SNE (top) and UMAP (bottom) visualisations of
         penultimate-layer embeddings for 5\,000 validation samples.}
\label{fig:tsne}
\end{figure}

The correlation between embedding quality and classification accuracy
(ConvNeXt $>$ ViT $>$ CNN in both metrics) suggests that ConvNeXt's
advantage stems from fundamentally better class-discriminative
representations rather than merely a better decision boundary.

% -----------------------------------------------------------------
\subsection{Data Cleaning Impact}\label{sec:cleaning}

\subsubsection{Cleaning Methodology}

Mislabelled and duplicate images were identified through a three-model
consensus pipeline: samples flagged as misclassified by all three
architectures (or detected as near-duplicates via perceptual hashing)
were surfaced for manual review.  This process identified 650~images
for removal and 1\,296~images for relabelling, yielding a cleaned
dataset of 203\,655~images (down from 204\,305).  The cleaned data was
re-split into 75/15/10 train/validation/test partitions
(152\,657\,/\,30\,511\,/\,20\,487 images).

\subsubsection{Retraining Results}

ConvNeXt-Base was retrained on the cleaned data with the same default
hyperparameters (lr$\,{=}\,$1e-4, wd$\,{=}\,$1e-4, batch$\,{=}\,$64,
10~epochs).  Table~\ref{tab:cleaning} compares before- and
after-cleaning performance.

\begin{table}[ht]
\centering
\caption{Data cleaning impact on ConvNeXt-Base.  The ``Before'' model
         was trained on the original 75/25 split; the ``After'' model on
         the cleaned 75/15/10 split with identical hyperparameters.}
\label{tab:cleaning}
\begin{tabular}{lccc}
\toprule
\textbf{Condition} & \textbf{Val Acc (\%)} & \textbf{Test Acc (\%)} & \textbf{Split} \\
\midrule
Before cleaning & 83.45 & ---\textsuperscript{a} & 75/25 \\
After cleaning  & 84.55 & 84.54 & 75/15/10 \\
\midrule
\textbf{Improvement} & \textbf{+1.10 pp} & --- & --- \\
\bottomrule
\multicolumn{4}{l}{\footnotesize\textsuperscript{a}No held-out test set
  existed in the original split.}
\end{tabular}
\end{table}

The $+1.10$~pp improvement from cleaning alone---without any
hyperparameter changes---validates the investment in data quality.
Test accuracy (84.54\%) closely matches validation accuracy (84.55\%),
indicating that the model generalises well and that the new
train/val/test split does not exhibit overfitting artefacts.

Residual hard classes persist after cleaning (Table~\ref{tab:hard-classes}).
The worst-performing classes on the test set are \emph{whole wheat}
(22.5\%), \emph{Bee hoon, soto} (32.3\%), and \emph{Mee siam, fried}
(44.4\%).  These classes suffer from inherent visual ambiguity with
near-identical counterparts (\emph{whole grain bread}, \emph{Mee soto},
\emph{Bee hoon, goreng}) rather than from labelling noise, suggesting
that further improvement would require taxonomy revision.

\begin{table}[ht]
\centering
\caption{Bottom-10 classes by test accuracy after cleaning
         (ConvNeXt-Base, cleaned data).}
\label{tab:hard-classes}
\begin{tabular}{lcc}
\toprule
\textbf{Class} & \textbf{Test Acc (\%)} & \textbf{Test Images} \\
\midrule
whole wheat             & 22.5 &  40 \\
Bee hoon, soto          & 32.3 &  62 \\
Mee siam, fried         & 44.4 &  45 \\
bee hoon                & 45.3 &  64 \\
Barley                  & 47.5 &  40 \\
tortilla -- plain       & 47.8 &  46 \\
Mee kuah                & 51.0 &  51 \\
chicken soup            & 57.4 & 136 \\
Don, chicken teriyaki   & 57.4 &  61 \\
Lontong w/ sayur lodeh  & 58.3 &  48 \\
\bottomrule
\end{tabular}
\end{table}

\textbf{Note:} ResNet-50 and ViT-Base were not individually retrained
on the cleaned data; the cleaning impact was measured exclusively
through ConvNeXt-Base.

% -----------------------------------------------------------------
\subsection{Learning Rate Sensitivity Analysis}\label{sec:lr-sensitivity}

To disentangle optimisation effects from architectural advantages, all
three models were trained for 5~epochs on the cleaned data at five
learning rates spanning two orders of magnitude:
$\{$1e-5, 5e-5, 1e-4, 3e-4, 1e-3$\}$.  All other hyperparameters were
held constant (weight decay$\,{=}\,$1e-4, batch size$\,{=}\,$32,
cosine-annealing schedule, basic augmentation).
Table~\ref{tab:lr-sweep} reports best validation accuracy at each
learning rate.

\begin{table}[ht]
\centering
\caption{Learning rate sensitivity: best validation accuracy (\%) after
         5~epochs on cleaned data (batch size$\,{=}\,$32, wd$\,{=}\,$1e-4).
         Bold indicates each model's best LR.}
\label{tab:lr-sweep}
\begin{tabular}{lccc}
\toprule
\textbf{Learning Rate} & \textbf{ResNet-50} & \textbf{ViT-Base} & \textbf{ConvNeXt-Base} \\
\midrule
1e-5 & 42.00 & 81.44 & 82.12 \\
5e-5 & 69.31 & \textbf{82.53} & 83.99 \\
1e-4 & 73.80 & 81.03 & \textbf{84.30} \\
3e-4 & 77.71 & 66.76 & 82.35 \\
1e-3 & \textbf{78.88} & 34.20 & 75.70 \\
\bottomrule
\end{tabular}
\end{table}

\begin{figure}[ht]
\centering
% \includegraphics[width=0.75\textwidth]{figures/lr_sensitivity_curves.pdf}
\textbf{[Placeholder: LR sensitivity curves --- validation accuracy
vs.\ learning rate for the three architectures.]}
\caption{Learning rate sensitivity profiles.  ConvNeXt-Base dominates
         across the full range; ViT-Base degrades sharply above 1e-4.}
\label{fig:lr-sensitivity}
\end{figure}

Three key findings emerge:

\begin{enumerate}
  \item \textbf{ConvNeXt-Base wins at every learning rate tested.}
    Even at ViT's optimal LR (5e-5), ConvNeXt achieves 83.99\% versus
    ViT's 82.53\%.  This confirms that the architecture gap is genuine
    and not an artefact of sub-optimal LR selection.

  \item \textbf{ViT-Base is highly LR-sensitive.}  Performance drops
    from 82.53\% at 5e-5 to 34.20\% at 1e-3---a catastrophic
    48.33~pp collapse.  This is consistent with the known instability
    of self-attention layers under large learning rates during
    fine-tuning, where large updates destabilise the pretrained
    attention patterns.

  \item \textbf{ResNet-50 is monotonically improving with higher LR.}
    Unlike the other two models, ResNet-50 does not exhibit a clear
    optimum within the tested range, peaking at the highest LR (1e-3)
    with 78.88\%.  This tolerance to high learning rates is
    characteristic of convolution-heavy architectures with inherently
    stable gradient flow through local receptive fields.
\end{enumerate}

The optimal learning rates identified are: ResNet-50 $=$ 1e-3,
ViT-Base $=$ 5e-5, and ConvNeXt-Base $=$ 1e-4.  These are used in the
final fair comparison (Section~\ref{sec:final-comparison}).

% -----------------------------------------------------------------
\subsection{Final Architecture Comparison at Optimal Learning Rates}
\label{sec:final-comparison}

\textit{This section reports the definitive comparison: each
architecture trained for 10~epochs at its individually optimal learning
rate (from Section~\ref{sec:lr-sensitivity}) with batch size$\,{=}\,$32,
weight decay$\,{=}\,$1e-4, and early stopping (patience$\,{=}\,$5) on
the cleaned 75/15/10 split.}

\begin{table}[ht]
\centering
\caption{Final architecture comparison at optimal learning rates
         (10~epochs, BS$\,{=}\,$32, cleaned data).}
\label{tab:final-comparison}
\begin{tabular}{lcccccc}
\toprule
\textbf{Model} & \textbf{Best LR} &
  \textbf{Val Acc (\%)} & \textbf{Test Acc (\%)} &
  \textbf{Top-5 Acc (\%)} & \textbf{Macro F1} & \textbf{Weighted F1} \\
\midrule
ResNet-50     & 1e-3 & \texttt{XX.XX} & \texttt{XX.XX} & \texttt{XX.XX} & \texttt{X.XXX} & \texttt{X.XXX} \\
ViT-Base      & 5e-5 & \texttt{XX.XX} & \texttt{XX.XX} & \texttt{XX.XX} & \texttt{X.XXX} & \texttt{X.XXX} \\
ConvNeXt-Base & 1e-4 & \texttt{XX.XX} & \texttt{XX.XX} & \texttt{XX.XX} & \texttt{X.XXX} & \texttt{X.XXX} \\
\bottomrule
\end{tabular}
\end{table}

\textbf{[TODO: Fill in Table~\ref{tab:final-comparison} after completing
Section~D of Notebook~12.  The 10-epoch training runs were interrupted
and need to be re-executed.]}

% -----------------------------------------------------------------
% SECTION 5 — DISCUSSION
% -----------------------------------------------------------------
\section{Discussion}\label{sec:discussion}

\subsection{Why ConvNeXt Outperforms}\label{sec:why-convnext}

ConvNeXt~\cite{liu2022convnet} modernises the standard CNN architecture
by incorporating design elements pioneered in Vision Transformers:
$7{\times}7$ depthwise convolutions that approximate the large receptive
fields of self-attention, LayerNorm instead of BatchNorm, GELU
activations, and an inverted bottleneck block structure.  This design
retains the CNN's strong inductive biases---translation equivariance and
locality---while gaining the ability to capture longer-range
dependencies.

Our results provide converging evidence for this advantage across
multiple evaluation axes:

\begin{itemize}
  \item \textbf{Accuracy:} ConvNeXt achieves the highest validation
    accuracy at every learning rate tested
    (Section~\ref{sec:lr-sensitivity}), ruling out optimisation-luck
    as an explanation.
  \item \textbf{Attention patterns:} Grad-CAM analysis
    (Section~\ref{sec:gradcam}) shows ConvNeXt balances local texture
    and global structure, unlike ResNet-50 (texture-biased) or ViT
    (structure-biased).
  \item \textbf{Embedding quality:} ConvNeXt produces the most
    separable embeddings (silhouette score 0.200 vs.\ 0.134/0.129),
    indicating fundamentally better learned representations
    (Section~\ref{sec:embeddings}).
  \item \textbf{Error patterns:} ConvNeXt exhibits fewer
    architecture-specific confusions (Section~\ref{sec:confusion-pairs}),
    with its remaining errors concentrated on inherently ambiguous class
    pairs rather than systematic feature blind-spots.
\end{itemize}

For Singaporean food classification in particular, where dishes
frequently share base ingredients (noodles, rice) but differ in
preparation style, sauce, or garnish, the ability to integrate both
textural detail and compositional context appears decisive.

\subsection{Dataset-Specific Challenges}\label{sec:dataset-challenges}

SGFood-233 presents several challenges beyond those encountered in
standard food benchmarks such as Food-101:

\begin{enumerate}
  \item \textbf{Inherently ambiguous classes:} Pairs such as
    \emph{char siew pau}~$\leftrightarrow$~\emph{steamed buns} and
    \emph{Lontong with sayur lodeh}~$\leftrightarrow$~\emph{lontong}
    are visually near-identical; accurate classification requires
    recognising subtle contextual cues (filling colour, side dishes).
  \item \textbf{Fine-grained noodle taxonomy:} Multiple \emph{bee hoon},
    \emph{mee}, and \emph{mee siam} variants share the same noodle base
    and differ primarily in sauce or broth, demanding
    context-aware feature extraction beyond local texture.
  \item \textbf{Residual labelling ambiguity:} Despite our cleaning
    effort, Tier~2 mislabels (borderline cases where even human
    annotators disagree) remain, particularly in soup and noodle
    categories.  Classes like \emph{whole wheat} (22.5\% test accuracy)
    and \emph{Bee hoon, soto} (32.3\%) likely require taxonomy
    consolidation rather than further model optimisation.
\end{enumerate}

\subsection{Impact of Data Cleaning}\label{sec:discuss-cleaning}

The $+1.10$~pp improvement from data cleaning
(Section~\ref{sec:cleaning}) demonstrates that label noise is a
meaningful bottleneck even for large-scale food datasets.  The
three-model consensus approach---flagging samples misclassified by
all architectures---proved effective and scalable, as it leverages the
complementary error profiles documented in
Section~\ref{sec:confusion-pairs} to surface true labelling errors
while reducing false positives.

Notably, the hardest classes after cleaning (\emph{whole wheat},
\emph{Bee hoon, soto}) are not noise-dominated but
taxonomy-ambiguous.  This suggests that further accuracy gains require
not additional cleaning but rather class hierarchy design---for example,
merging \emph{whole wheat} and \emph{whole grain bread} into a single
category or introducing a hierarchical label structure for noodle
subcategories.

\subsection{Learning Rate Sensitivity Insights}\label{sec:discuss-lr}

The LR sweep (Section~\ref{sec:lr-sensitivity}) reveals that learning
rate is the single most impactful hyperparameter when fine-tuning
pretrained models on a new food dataset, with performance varying by up
to 48~pp for ViT-Base.

Architecture determines LR tolerance: convolution-heavy models
(ResNet-50, ConvNeXt) tolerate wider LR ranges due to stable gradient
flow through local receptive fields, while the fully self-attention-based
ViT requires careful tuning.  ConvNeXt occupies a favourable middle
ground---it peaks at the commonly-used 1e-4 and degrades gracefully at
both extremes---making it the easiest architecture to tune in practice.

This has a direct practical implication: for practitioners applying
transfer learning to new food classification tasks, ConvNeXt's LR
robustness reduces the risk of sub-optimal hyperparameter selection,
lowering the barrier to achieving strong performance.

\subsection{Practical Deployment Considerations}\label{sec:deployment}

Table~\ref{tab:deployment} maps the experimental results to deployment
scenarios.

\begin{table}[ht]
\centering
\caption{Architecture recommendations by deployment scenario.}
\label{tab:deployment}
\begin{tabular}{lll}
\toprule
\textbf{Scenario} & \textbf{Recommended Model} & \textbf{Rationale} \\
\midrule
Maximum accuracy     & ConvNeXt-Base & Best accuracy and embeddings \\
Mobile / edge        & ResNet-50     & 3.7$\times$ fewer params, 1.8$\times$ faster \\
Balanced production  & ConvNeXt-Base & Strong accuracy with acceptable throughput \\
\bottomrule
\end{tabular}
\end{table}

For mobile deployment where ResNet-50's $-8.36$~pp accuracy penalty is
unacceptable, a smaller ConvNeXt variant (e.g.\ ConvNeXt-Tiny, 28.6\,M
parameters) or knowledge distillation from ConvNeXt-Base into a compact
student model would be promising directions.

\subsection{Limitations}\label{sec:limitations}

This study has several limitations that should be considered when
interpreting the results:

\begin{enumerate}
  \item \textbf{Single dataset.}  All experiments were conducted on
    SGFood-233; generalisability to other food corpora (e.g.\ Food-101,
    ISIA Food-500) is not established.
  \item \textbf{Fixed input resolution.}  All models used
    $224{\times}224$-pixel inputs.  Higher resolutions may
    disproportionately benefit architectures that rely on fine-grained
    spatial detail.
  \item \textbf{Class imbalance.}  Class frequencies in SGFood-233
    range from $\sim$200 to $\sim$5\,000 images.  No explicit
    re-sampling or loss re-weighting was applied; low-frequency classes
    may be under-represented in aggregate metrics.
  \item \textbf{Hyperparameter search scope.}  Only learning rate was
    systematically varied; weight decay, augmentation strategy, and
    label smoothing were held at defaults.  A broader search
    (e.g.\ Bayesian optimisation) might narrow or widen the
    architecture gaps observed.
  \item \textbf{Single seed.}  All experiments used seed$\,{=}\,$42.
    Without multiple runs, we cannot report confidence intervals or
    assess statistical significance of the observed accuracy
    differences.
  \item \textbf{No cross-architecture cleaning comparison.}  The data
    cleaning impact was measured by retraining ConvNeXt-Base only;
    ResNet-50 and ViT-Base may benefit differently from the same
    cleaning.
\end{enumerate}

\subsection{Future Work}\label{sec:future-work}

Several directions could extend and strengthen these findings:

\begin{itemize}
  \item \textbf{Comprehensive hyperparameter optimisation:}  Extend the
    search to weight decay, augmentation policy (RandAugment, CutMix),
    and label smoothing for ConvNeXt-Base.
  \item \textbf{Cross-dataset validation:}  Evaluate on Food-101 and
    ISIA Food-500 to test whether ConvNeXt's advantage generalises
    beyond Singaporean food.
  \item \textbf{Knowledge distillation:}  Distil ConvNeXt-Base into a
    mobile-friendly student (e.g.\ MobileNetV3) to combine strong
    accuracy with edge-device constraints.
  \item \textbf{Hierarchical classification:}  Introduce a two-level
    taxonomy (e.g.\ \emph{noodle}~$\to$~\emph{bee hoon goreng}) to
    reduce confusion among fine-grained subcategories.
  \item \textbf{Higher input resolution:}  Test $384{\times}384$ inputs,
    which may particularly benefit noodle and soup classes where
    distinguishing details are small.
\end{itemize}

% -----------------------------------------------------------------
% SECTION 6 — CONCLUSION
% -----------------------------------------------------------------
\section{Conclusion}\label{sec:conclusion}

This study presented a systematic comparison of three transfer-learning
architectures---ResNet-50 (CNN), ViT-Base (Transformer), and
ConvNeXt-Base (modernised CNN)---for classifying 233 categories of
Singaporean food in the SGFood-233 dataset (${\sim}$203\,K images).

ConvNeXt-Base consistently emerged as the strongest architecture across
every evaluation axis.  Under identical default hyperparameters, it
achieved 83.45\% validation accuracy, outperforming ResNet-50 by
$+8.36$~pp and ViT-Base by $+3.29$~pp.  A learning rate sensitivity
analysis across five rates spanning two orders of magnitude confirmed
that this advantage is architectural rather than optimisation-dependent:
ConvNeXt-Base ranked first at \emph{every} learning rate tested,
reaching 84.30\% at its optimal rate (1e-4) after just 5~epochs on
cleaned data.  Complementary analyses showed that ConvNeXt-Base
produces the most class-separable embeddings (silhouette
score~$=$~0.200 vs.\ 0.134/0.129), balances local and global attention
in Grad-CAM maps, and exhibits fewer systematic confusion patterns.

A data cleaning pipeline based on three-model consensus flagging
identified and corrected 1\,946 problematic images (650 removed, 1\,296
relabelled).  Retraining ConvNeXt-Base on the cleaned data improved
validation accuracy from 83.45\% to 84.55\% ($+1.10$~pp) with a
matching test accuracy of 84.54\%, demonstrating that targeted data
quality investment yields measurable gains even for large-scale
datasets.

% Uncomment and update the line below after completing Section D:
% The definitive 10-epoch comparison at individually optimal learning
% rates yielded \texttt{XX.XX}\% / \texttt{XX.XX}\% / \texttt{XX.XX}\%
% validation accuracy for ResNet-50 / ViT-Base / ConvNeXt-Base
% respectively, further confirming ConvNeXt-Base as the recommended
% backbone for Singaporean food image classification.

Based on these findings, we recommend ConvNeXt-Base as the default
backbone for Singaporean food image classification, with ResNet-50 as a
viable alternative for latency-critical mobile deployments.  Future work
should expand the hyperparameter search beyond learning rate, validate on
additional food datasets, and explore knowledge distillation to bridge
the accuracy--efficiency gap for on-device inference.
